%% ===================================================================
\section{Experimental Results}
\label{sec:eval}

\subsection{PolyBench/C (30 Kernels)}

\begin{table}[t]
  \caption{PolyBench/C results at 1$\times$ sizes. Agent = autonomous LLM with tool use (no analysis); NA = LLM with classified error feedback (no analysis); WA+Prof = our method (compiler analysis + NCU profiling feedback).}
  \label{tab:results_1x}
  \centering\small
  \begin{tabular}{lccc}
    \toprule
    \textbf{Metric} & \textbf{Agent} & \textbf{NA} & \textbf{WA+Prof} \\
    \midrule
    Pass rate          & 24/30 (80\%) & 28/30 (93\%) & 29/30 (97\%) \\
    First-try passes   & 11           & 8            & 12           \\
    Median speedup     & 0.60$\times$ & 0.88$\times$ & 0.76$\times$ \\
    Mean speedup       & 0.87$\times$ & 1.64$\times$ & 1.24$\times$ \\
    Kernels $>1\times$ & 7/23         & 12/28        & 11/29        \\
    \bottomrule
  \end{tabular}
\end{table}

\begin{table}[t]
  \caption{PolyBench/C results at 8$\times$ sizes ($N$: 60$\to$480, matrices: 120$\to$960). Speedups vs.\ single-threaded C.}
  \label{tab:results_8x}
  \centering\small
  \begin{tabular}{lccc}
    \toprule
    \textbf{Metric} & \textbf{Agent} & \textbf{NA} & \textbf{WA+Prof} \\
    \midrule
    Pass rate          & 17/30 (57\%) & 26/30 (87\%) & 27/30 (90\%) \\
    Median speedup     & 20.7$\times$ & 10.8$\times$ & 24.6$\times$ \\
    Kernels $>1\times$ & 13/15        & 14/20        & 18/21        \\
    \bottomrule
  \end{tabular}
\end{table}

Tables~\ref{tab:results_1x} and~\ref{tab:results_8x} compare three configurations on PolyBench/C.
Our method (WA+Prof) achieves the highest correctness at both scales: 97\% at 1$\times$ and 90\% at 8$\times$.
The autonomous agent baseline, despite having tool access (write code, run tests, run benchmarks), achieves only 80\% at 1$\times$ and 57\% at 8$\times$---the weakest of all three.
This shows that autonomous tool use without compiler analysis guidance leads to poor retry strategies: the agent wastes turns on ineffective fixes without understanding the root cause of parallelization failures.

At 1$\times$ sizes, median speedups are limited by GPU launch overhead at small $N$.
At 8$\times$ sizes, WA+Prof achieves 24.6$\times$ median, outperforming NA (10.8$\times$) and the agent (20.7$\times$).
Top kernels: \texttt{3mm} (2041$\times$), \texttt{2mm} (2004$\times$), \texttt{gemm} (451$\times$).

\subsection{Profiling-Guided Optimization}

\begin{table}[t]
  \caption{Representative profiling feedback improvements.
    ``Before'' = first correct kernel; ``After'' = best after NCU-guided iterations.}
  \label{tab:profiling}
  \centering\small
  \begin{tabular}{lccl}
    \toprule
    \textbf{Kernel} & \textbf{Before} & \textbf{After} & \textbf{Scale} \\
    \midrule
    3mm         & 0.10$\times$ & 4.01$\times$ & 1$\times$ \\
    covariance  & 0.11$\times$ & 3.66$\times$ & 1$\times$ \\
    2mm         & 1.48$\times$ & 2.97$\times$ & 1$\times$ \\
    syrk        & 1.65$\times$ & 2.07$\times$ & 1$\times$ \\
    \midrule
    gemm        & 120$\times$  & 451$\times$  & 8$\times$ \\
    correlation & 12$\times$   & 51$\times$   & 8$\times$ \\
    atax        & 0.4$\times$  & 15$\times$   & 8$\times$ \\
    syrk        & 24$\times$   & 62$\times$   & 8$\times$ \\
    \bottomrule
  \end{tabular}
\end{table}

Profiling feedback improved 9 of 29 kernels at 1$\times$ and 12 of 21 benchmarked kernels at 8$\times$ (Table~\ref{tab:profiling}).
The gains are largest on latency-bound kernels where the initial implementation has insufficient parallelism: the NCU metrics reveal low SM and memory utilization, and the LLM restructures the kernel accordingly.
Correctness is preserved: each optimized kernel is re-validated, and attempts that break correctness are rejected.

\subsection{TSVC (151 Kernels)}

\sysname passes 143/151 TSVC kernels (95\%), 91 on first attempt,
with a mean speedup of 2.78$\times$ and maximum of 267$\times$ over single-threaded C.

\subsection{Application Kernels}

\begin{table}[t]
  \caption{Application kernels: Triton speedup vs.\ single-threaded C.}
  \label{tab:realworld}
  \centering\small
  \begin{tabular}{llcc}
    \toprule
    \textbf{Source} & \textbf{Kernel} & \textbf{WA/ST} & \textbf{Pass?} \\
    \midrule
    ECP     & LJ force   & 24.1$\times$ & Y \\
    ECP     & SpMV       & 7.1$\times$  & Y \\
    Rodinia & SRAD       & 12.8$\times$ & Y \\
    Rodinia & pathfinder & 3.6$\times$  & Y \\
    Rodinia & lud        & 1.3$\times$  & Y \\
    Rodinia & gaussian   & 0.6$\times$  & Y \\
    \midrule
    Rodinia & hotspot    & ---          & N \\
    Rodinia & lavaMD     & ---          & N \\
    \bottomrule
  \end{tabular}
\end{table}

We apply \sysname to 3 Rodinia and 5 ECP/Rodinia application kernels (Table~\ref{tab:realworld}).
6 of 8 pass correctness.
The highest speedups come from compute-intensive kernels:
\texttt{LJ\_force} (24.1$\times$) and \texttt{SRAD} (12.8$\times$).
Two kernels fail: \texttt{hotspot} (complex 2D stencil with boundary conditions) and \texttt{lavaMD} (irregular particle interactions), both requiring domain-specific optimizations beyond what the unified analysis provides.
